% DSAIT4335 Recommender Systems -- Final Project report (hand-in version)
%
% Formatting required by the brief: Times New Roman, 12 pt, line spacing 1.15.
% Structure required by the brief: first page = group information; main body = at most ONE page per task
% (results + discussion) with at most 200 WORDS of discussion per task; appendix = additional results
% (same limits). The main body must be self-contained.
%
% Who writes what (RecSys-Project-Plan-and-Team-Split, Part 2): each member writes the subsections of the
% sub-tasks they own; A edits the whole report for one voice and the page limits.
%   A  Models        1.1, 1.2, 2.2        C  Weighted hybrid  1.3, 1.5 (weighted), 2.3, 2.6
%   B  Evaluation    2.1, 2.5, 3.4        D  Content-based    1.1 (content-based), 1.4, 1.5 (others), 2.4, 3.3
%   E  Re-rankers    3.1, 3.2
%
% Build: cd report && tectonic report.tex      (or: latexmk -pdf -outdir=build report.tex)
% Overleaf: Menu -> Main document -> report/report.tex
%
% Generated assets (committed to git, produced by python -m project.experiments.<...>):
%   ../figures/generated/<name>.pdf        -> \generatedfigure[width=...]{<name>}
%   tables/generated/<name>.tex            -> \generatedtable{<name>}
% A red box is shown for an asset that has not been generated/committed yet, so the report always compiles.
\documentclass[12pt,a4paper]{article}

\usepackage[T1]{fontenc}
\usepackage[utf8]{inputenc}
\usepackage{newtxtext,newtxmath}          % Times New Roman (text + maths)
\usepackage{setspace}
\setstretch{1.15}                         % line spacing 1.15
\usepackage[a4paper,margin=2.2cm]{geometry}
\usepackage{graphicx}
\usepackage{booktabs}
\usepackage{amsmath,amssymb}
\usepackage{xcolor}
\usepackage[font=small,skip=4pt]{caption}
\usepackage{enumitem}
\setlist{nosep,leftmargin=*}
\usepackage[hidelinks]{hyperref}
\usepackage{titlesec}
\titlespacing*{\section}{0pt}{8pt}{4pt}
\titlespacing*{\subsection}{0pt}{6pt}{2pt}

% ---------------------------------------------------------------------------
% Title-page data -- fill in
% ---------------------------------------------------------------------------
\newcommand{\coursecode}{DSAIT4335}
\newcommand{\coursename}{Recommender Systems}
\newcommand{\projecttitle}{Hybrid Recommendation on MovieLens-100K:\\ Accuracy, Coefficients and Societal Re-ranking}
\newcommand{\groupnumber}{\TODO{group number}}

% ---------------------------------------------------------------------------
% Helpers
% ---------------------------------------------------------------------------
\newcommand{\TODO}[1]{\textcolor{red}{\textbf{[TODO: #1]}}}
\newcommand{\owner}[1]{\marginpar{\footnotesize\color{gray}#1}}   % who writes this subsection (remove before hand-in)

\newcommand{\missingasset}[1]{%
  \fbox{\parbox{0.85\linewidth}{\centering\color{red}%
    \textbf{TODO:} generated asset not found\\ \texttt{#1}}}}

\graphicspath{{../figures/generated/}{figures/}}
\newcommand{\generatedfigure}[2][width=\linewidth]{%
  \IfFileExists{../figures/generated/#2.pdf}%
    {\includegraphics[#1]{../figures/generated/#2.pdf}}%
    {\IfFileExists{figures/generated/#2.pdf}%
      {\includegraphics[#1]{figures/generated/#2.pdf}}%
      {\missingasset{figures/generated/\detokenize{#2}.pdf}}}}
\newcommand{\generatedtable}[1]{%
  \IfFileExists{tables/generated/#1.tex}%
    {\input{tables/generated/#1.tex}}%
    {\IfFileExists{report/tables/generated/#1.tex}%
      {\input{report/tables/generated/#1.tex}}%
      {\missingasset{report/tables/generated/\detokenize{#1}.tex}}}}

\begin{document}

% ===========================================================================
% First page: group information
% ===========================================================================
\pagenumbering{roman}
\begin{titlepage}
  \centering
  \vspace*{3cm}
  {\large \coursecode\ \coursename\ --- Final Project\par}
  \vspace{1.5cm}
  {\LARGE\bfseries \projecttitle\par}
  \vspace{1.5cm}
  {\Large Group \groupnumber\par}
  \vspace{1.5cm}
  {\large\bfseries Group members\par}
  \vspace{0.5em}
  \begin{tabular}{ll}
    \TODO{name} & \TODO{student number}\\
    \TODO{name} & \TODO{student number}\\
    \TODO{name} & \TODO{student number}\\
    \TODO{name} & \TODO{student number}\\
    \TODO{name} & \TODO{student number}\\
  \end{tabular}
  \vfill
  {\large Delft University of Technology\par}
  {\normalsize Code: \texttt{project/} in the submitted archive (see \texttt{project/models/README.md} for how to run it)\par}
\end{titlepage}
\pagenumbering{arabic}

% ===========================================================================
% TASK 1 -- Hybrid recommender (max. one page, max. 200 words of discussion)
% ===========================================================================
\section{Task 1 --- Hybrid Recommender}
\label{sec:task1}

\subsection{Individual recommender models (Tasks 1.1 and 1.2)}\owner{A}
\label{sec:task1-individual}
\paragraph{Setup (shared by all tasks).}
All experiments use RecBole on MovieLens-100K with \emph{one} frozen split: seed 2020, a random
80/10/10 split per user (RecBole \texttt{RS}, \texttt{RO}, \texttt{group\_by: user}), full ranking over all
items the user has not interacted with, and every held-out interaction counted as relevant (RecBole's default; no
rating threshold). Train/validation/test hold 80\,808 / 9\,596 / 9\,596 interactions of 943 users and 1\,682
items; the split is exported once with a checksum manifest and verified to be identical through every model's
configuration. Every model exports its score for all user--item pairs and a top-50 candidate list per user
(validation lists mask the training history, test lists mask train\,+\,validation); the final list is the top 10.
Hybrids and re-rankers only read these files, so nothing is retrained downstream.

\paragraph{Individual models (Task 1.1).}
We use the course's RecBole models, one per family discussed in class, plus the two na\"ive baselines: Random
(scores drawn per user; RecBole's own \texttt{Random} gives all users of a batch the same list) and Pop;
neighbourhood models UserKNN and ItemKNN (cosine, RecBole \texttt{ItemKNN} with \texttt{knn\_method});
the linear item--item models EASE and SLIMElastic; BPR matrix factorisation (pairwise ranking loss);
the neural NeuMF and FISM; and the graph models LightGCN and NGCF. The content-based model of
Assignment~1 is added by Track~D on the same split.

\paragraph{Tuning (Task 1.2).}
Each model is tuned by an exhaustive grid (7--20 configurations, Table~\ref{tab:tuning-summary} in the appendix)
on \emph{validation} NDCG@10, the target agreed for every later choice in the project; the test set is used
once, for the selected configuration. Neural models train for up to 100 epochs with validation after every
epoch and early stopping after 10 epochs without improvement, instead of the course configs' fixed 20 epochs.
Selected: ItemKNN $k{=}200$, UserKNN $k{=}50$ (both without shrinkage), EASE $\lambda{=}500$, SLIM
$\alpha{=}0.1$, $l_1{=}0.01$, BPR 128 factors with learning rate $10^{-3}$, NeuMF MLP $[128,64,32]$ with
dropout 0.1, LightGCN 3 layers with learning rate $5\cdot10^{-3}$.

\emph{Discussion.} Tuning matters most where the course configuration left a model under-trained: LightGCN
rises from 0.136 to 0.243 validation NDCG@10 (its best epoch is 94 of 100) and BPR from 0.199 to 0.241, both
mainly through a higher learning rate or a longer schedule; the memory-based and linear models move by at most
0.003 because, with 85 interactions per user on average, item--item statistics are already well estimated
and extra neighbours or shrinkage barely change the ranking. EASE prefers strong regularisation
($\lambda{=}500$; $\lambda{=}10$ loses 0.05), consistent with a dense, small catalogue in which an
unregularised item--item solution overfits co-occurrence noise.

\TODO{D: one sentence on the content-based model (Assignment 1 BERT model on the same split, exported as a score file).}

\subsection{Weighted hybrid with learned coefficients (Task 1.3)}\owner{C}
\TODO{C: regression model used to learn the coefficients; fitted on the validation candidates (union of each model's top-50); per-user score normalisation; the learned weights.}

\subsection{Other hybrid models (Tasks 1.4 and 1.5)}\owner{D, C}
\TODO{D: switching and mixed hybrids and why; C/D: how all hybrids were tuned on validation NDCG@10.}

\clearpage
% ===========================================================================
% TASK 2 -- Evaluation of effectiveness (max. one page, max. 200 words of discussion)
% ===========================================================================
\section{Task 2 --- Evaluation of Effectiveness}
\label{sec:task2}

\subsection{Evaluation metrics (Task 2.1)}\owner{B}
\TODO{B: accuracy and beyond-accuracy metrics as implemented in \texttt{project/metrics} (validated against RecBole on the Track A lists).}

\subsection{Comparison with the baselines (Task 2.2)}\owner{A}
\label{sec:task2-baselines}
Table~\ref{tab:model-results} reports the tuned models on the test split (accuracy metrics computed by
RecBole on the exported lists; beyond-accuracy columns are added from Section~\ref{sec:task2} once Track~B's
metrics are in; FISM and NGCF are reported with the course configuration, untuned).

\begin{table}[htbp]
  \centering
  \caption{Test-set accuracy of the individual models (lists of 10, all held-out interactions relevant); best
  value per column in bold. The last column is the validation NDCG@10 used for model selection.}
  \label{tab:model-results}
  \small\setlength{\tabcolsep}{4pt}
  \generatedtable{model_results}
\end{table}

\emph{Discussion.} Pop is a strong na\"ive baseline here: one unpersonalised list hits a relevant film for
47\% of the users (Random: 7\%) because MovieLens-100K's interactions concentrate on a few hundred films.
All personalised models beat it, yet 47--61\% of their top-10 slots are films Pop also lists, and their
recommendations were rated by 23--27\% of the training users on average (Pop 33\%, Random under 5\%).
The item--item linear models EASE and SLIM are the most accurate: with 943 users and 1\,682 items the
co-occurrence matrix is dense enough for a closed-form, strongly regularised solution, whereas the
embedding models need long training to reach a similar level on this little data. The accuracy winners
thus inherit much of Pop's popularity bias, which Tasks~2.5 and~3 examine.

\subsection{Hybrid coefficient analysis (Task 2.3)}\owner{C}
\TODO{C: contribution of each individual recommender to the hybrid (coefficients, drop-one ablation).}

\subsection{Further analyses (Task 2.4)}\owner{D}
\TODO{D: overlap between the models' lists; what an oracle best-model-per-user would score.}

\subsection{User and item groups (Task 2.5)}\owner{B}
\TODO{B: accuracy and beyond-accuracy per user group (activity, popularity taste) and item group (head/mid/tail).}

\subsection{Insights (Task 2.6)}\owner{C}
\TODO{C: group-specific hybrid motivated by 2.5.}

\clearpage
% ===========================================================================
% TASK 3 -- Societal aspects (max. one page, max. 200 words of discussion)
% ===========================================================================
\section{Task 3 --- Societal Aspects}
\label{sec:task3}

\subsection{Re-rankers (Task 3.1)}\owner{E}
\TODO{E: diversity (MMR), calibration, user-side and item-side fairness re-rankers on the top-50 candidates.}

\subsection{Effect of the re-rankers (Task 3.2)}\owner{E}
\TODO{E: trade-off curves (NDCG@10 against the target metric over $\lambda$) on the best models and the hybrid.}

\begin{table}[htbp]
  \centering
  \caption{Results with re-rankers applied.}
  \label{tab:reranker-results}
  \generatedtable{reranker_results}
\end{table}

\subsection{Hybrids combined with re-rankers (Task 3.3)}\owner{D}
\TODO{D: re-rank then combine versus combine then re-rank.}

\subsection{Impact on user and item groups (Task 3.4)}\owner{B}
\TODO{B: before/after re-ranking per group.}

\clearpage
% ===========================================================================
% APPENDIX -- additional results (same limits per task)
% ===========================================================================
\appendix
\section{Additional results for Task 1}
\label{app:task1}
\subsection{Hyper-parameter search (Task 1.2)}
Search spaces (\texttt{project/configs/hyper/}): ItemKNN and UserKNN $k\in\{10,20,50,100,200\}$,
shrink $\in\{0,10,50,100\}$; EASE $\lambda\in\{10,50,100,250,500,1000,2000\}$; SLIMElastic
$\alpha\in\{0.1,0.2,0.5,1\}$, $l_1\in\{0.01,0.02,0.05,0.1\}$; BPR factors $\in\{32,64,128\}$,
learning rate $\in\{5\cdot10^{-4},10^{-3},5\cdot10^{-3}\}$; NeuMF learning rate $\in\{10^{-3},5\cdot10^{-3}\}$,
MLP $\in\{[64,32],[128,64,32]\}$, dropout $\in\{0.1,0.3\}$; LightGCN layers $\in\{1,2,3\}$, learning rate
$\in\{10^{-3},5\cdot10^{-3}\}$, $L_2\in\{10^{-4},10^{-3}\}$. Every configuration is listed in
\texttt{results/processed/tuning/<Model>.csv}. Two selected values sit on the border of their grid
(ItemKNN $k{=}200$, BPR 128 factors); the neighbouring value is within 0.003 NDCG@10, so the grids were not
extended. All numbers are single runs with seed 2020; a 3-seed variance check is left for the final release.

\begin{table}[htbp]
  \centering
  \caption{Grid search per model: number of configurations, selected values, and NDCG@10 before (course
  configuration) and after tuning on validation and test.}
  \label{tab:tuning-summary}
  \small\setlength{\tabcolsep}{4pt}
  \generatedtable{tuning_summary}
\end{table}

\subsection{Accuracy overview (Task 2.2)}
\begin{figure}[htbp]
  \centering
  \generatedfigure[width=\linewidth]{model_comparison}
  \caption{Test-set NDCG@10, Recall@10 and Precision@10 per model (hatched: baselines; dotted line: Pop).}
  \label{fig:model-comparison}
\end{figure}

\section{Additional results for Task 2}
\TODO{B, C, D: full per-group tables, extra metrics, alternative relevance rule.}

\section{Additional results for Task 3}
\TODO{E: extra $\lambda$ values, per-group tables after re-ranking.}

\end{document}
